\chapter{MARCO TEORICO}
El presente capítulo expone los fundamentos conceptuales que sustentan el desarrollo de una aplicación móvil orientada a fomentar micro interacciones sociales positivas entre jóvenes. Para ello…

\section{INTERACCIÓN HUMANO-COMPUTADORA (IHC)}
La Interacción Humano-Computadora es una disciplina multidisciplinar que estudia el diseño, la evaluación y la implementación de sistemas informáticos interactivos para el uso humano. El grupo de interés especial SIGCHI de la Asociación para la Maquinaria Informática, la define como la "disciplina relacionada con el diseño, la evaluación e implementación de sistemas informáticos interactivos para el uso de los seres humanos, y con el estudio de los fenómenos más importantes con los que está relacionado" (Hernández Bieliukas, 2014, p. 13). Esta disciplina no se limita al hardware o el software de manera aislada; abarca también los modelos mentales de los usuarios, las tareas que desempeña el sistema, su adaptación a distintos contextos y el impacto organizacional y social de la tecnología interactiva.

Según Diaper (1989), citado por Hernández Bieliukas (2014), los objetivos centrales de la IHC son desarrollar o mejorar la seguridad, utilidad, efectividad, eficiencia y usabilidad de los sistemas que incluyan computadoras. Para lograrlo, Preece (1994), referenciado en el mismo texto, plantea que es necesario comprender los factores psicológicos, ergonómicos, organizativos y sociales que determinan cómo las personas trabajan con los ordenadores, para trasladar ese conocimiento al diseño de herramientas que permitan una interacción eficiente, efectiva y segura, tanto a nivel individual como grupal (p. 14). Esta perspectiva establece que la IHC no puede reducirse a un problema técnico, sino que es fundamentalmente un problema humano.

La evolución histórica de la IHC muestra un desplazamiento progresivo de enfoque. Hernández Bieliukas (2014) señala que la disciplina surgió alrededor de 1970, cuando los computadores comenzaron a masificarse, pero los diseñadores asumían erróneamente que los usuarios tendrían un alto conocimiento técnico. La realidad demostró que los usuarios típicos se sentían frustrados con sistemas difíciles de usar, lo que impulsó el giro hacia la usabilidad como criterio central de diseño (p. 206). A mediados de los años noventa, con la aparición de los medios digitales, web, dispositivos móviles, televisión interactiva, se incorporó un nuevo componente a la IHC: las emociones, a partir de lo cual surgió el concepto de experiencia de usuario (p. 207).

